
\documentclass{article}
\usepackage{colortbl}
\usepackage{makecell}
\usepackage{multirow}
\usepackage{supertabular}

\begin{document}

\newcounter{utterance}

\centering \large Interaction Transcript for game `ifeval', experiment `ifeval', episode 488 with reflexion{-}llama{-}3.1{-}8b{-}instruct.
\vspace{24pt}

{ \footnotesize  \setcounter{utterance}{1}
\setlength{\tabcolsep}{0pt}
\begin{supertabular}{c@{$\;$}|p{.15\linewidth}@{}p{.15\linewidth}p{.15\linewidth}p{.15\linewidth}p{.15\linewidth}p{.15\linewidth}}
    \# & \multicolumn{2}{c}{Player} && \multicolumn{2}{c}{Game Master} \\
    \hline

    \theutterance \stepcounter{utterance}  
    & & & \multicolumn{4}{p{0.6\linewidth}}{
        \cellcolor[rgb]{0.9,0.9,0.9}{
            \makecell[{{p{\linewidth}}}]{
                \texttt{\tiny{[P1$\langle$GM]}}
                \texttt{You have attempted this same task before and did not succeed. These are the plans you wrote after those attempts.} \\
\texttt{Trial \#1: New plan: My previous attempt at writing a song about innovation was successful in terms of formatting, but it failed to capture the essence of the task, which was to write a song appealing to teenagers. I focused too much on the general theme of innovation and didn't specifically address the tone and language that would resonate with teenagers. Next time, I will focus on using more relatable language, catchy rhymes, and a more energetic tone to make the song more appealing to the target audience. I will also try to incorporate more contemporary references and examples that teenagers can relate to.} \\
\\ 
\texttt{Trial \#2: New plan: My previous attempt at writing a song about innovation was successful in terms of formatting, but it failed to capture the essence of the task, which was to write a song appealing to teenagers. I focused too much on the general theme of innovation and didn't specifically address the tone and language that would resonate with teenagers. However, upon re{-}examining my previous attempt, I realized that I did use some contemporary references and catchy rhymes, but I didn't go far enough in making the language relatable and energetic.} \\
\\ 
\texttt{To improve, I will focus on using more colloquial language, slang, and idioms that are commonly used by teenagers. I will also try to incorporate more specific examples and references to popular culture that are relevant to the teenage audience. Additionally, I will pay closer attention to the rhythm and flow of the song, using more consistent and catchy rhymes to make it more engaging and memorable.} \\
\\ 
\texttt{To achieve this, I will start by brainstorming a list of contemporary references, slang, and idioms that are commonly used by teenagers. I will then use this list to inform my writing, incorporating these elements into the song in a way that feels natural and authentic. I will also experiment with different rhythms and rhymes, using online tools and resources to help me find the right balance of flow and catchiness.} \\
\\ 
\texttt{Ultimately, my goal is to write a song that not only meets the technical requirements of the task but also resonates with the target audience of teenagers. By focusing on relatable language, catchy rhymes, and a more energetic tone, I am confident that I can create a song that is both engaging and memorable.} \\
\\ 
\texttt{Trial \#3: New plan: My previous attempts at writing a song about innovation were successful in terms of formatting, but they failed to capture the essence of the task, which was to write a song appealing to teenagers. I focused too much on the general theme of innovation and didn't specifically address the tone and language that would resonate with teenagers. However, upon re{-}examining my previous attempts, I realized that I did use some contemporary references and catchy rhymes, but I didn't go far enough in making the language relatable and energetic.} \\
\\ 
\texttt{To improve, I will focus on using more colloquial language, slang, and idioms that are commonly used by teenagers. I will also try to incorporate more specific examples and references to popular culture that are relevant to the teenage audience. Additionally, I will pay closer attention to the rhythm and flow of the song, using more consistent and catchy rhymes to make it more engaging and memorable.} \\
\\ 
\texttt{To achieve this, I will start by brainstorming a list of contemporary references, slang, and idioms that are commonly used by teenagers. I will then use this list to inform my writing, incorporating these elements into the song in a way that feels natural and authentic. I will also experiment with different rhythms and rhymes, using online tools and resources to help me find the right balance of flow and catchiness.} \\
\\ 
\texttt{However, I realize that my previous attempts were too focused on listing off examples of innovation and didn't really tell a story or convey a clear message. To improve, I will try to create a more narrative{-}driven song that tells a story about innovation and its impact on teenagers. I will use the contemporary references and slang to make the story feel more relatable and authentic, and I will focus on creating a clear and compelling message that resonates with the target audience.} \\
\\ 
\texttt{Specifically, I will try to write a song that addresses the following themes:} \\
\\ 
\texttt{* The excitement and possibility of innovation} \\
\texttt{* The importance of taking risks and trying new things} \\
\texttt{* The impact of innovation on teenagers' lives and futures} \\
\texttt{* The need for teenagers to be creative and think outside the box} \\
\\ 
\texttt{By focusing on these themes and using contemporary references and slang to make the language relatable and energetic, I am confident that I can create a song that is both engaging and memorable.} \\
\\ 
\\ 
\texttt{Write a song about innovation with a positive tone that is appealing to teenagers. Put your entire response in double quotation marks.} \\
            }
        }
    }
     \\ \\

    \theutterance \stepcounter{utterance}  
    & \multicolumn{4}{p{0.6\linewidth}}{
        \cellcolor[rgb]{0.561,0.737,0.561}{
            \makecell[{{p{\linewidth}}}]{
                \texttt{\tiny{[P1$\rangle$GM]}}
                \texttt{"Verse 1:} \\
\texttt{Yo, listen up, it's time to ignite,} \\
\texttt{The spark of innovation, it's feeling right,} \\
\texttt{A new idea's born, and it's taking flight,} \\
\texttt{The possibilities are endless, day and night.} \\
\\ 
\texttt{Chorus:} \\
\texttt{Innovate, don't hesitate,} \\
\texttt{Take the leap, don't be late,} \\
\texttt{Create, disrupt, and make it new,} \\
\texttt{That's what innovation's all about, boo.} \\
\\ 
\texttt{Verse 2:} \\
\texttt{From coding to art, to science and more,} \\
\texttt{There's a world of innovation waiting in store,} \\
\texttt{Don't be afraid to take a risk and try,} \\
\texttt{You never know what you might discover, you'll fly.} \\
\\ 
\texttt{Chorus:} \\
\texttt{Innovate, don't hesitate,} \\
\texttt{Take the leap, don't be late,} \\
\texttt{Create, disrupt, and make it new,} \\
\texttt{That's what innovation's all about, boo.} \\
\\ 
\texttt{Bridge:} \\
\texttt{It's not just about tech, it's about the mind,} \\
\texttt{Thinking outside the box, leaving the ordinary behind,} \\
\texttt{Be bold, be fearless, be you,} \\
\texttt{That's where the magic happens, it's true.} \\
\\ 
\texttt{Chorus:} \\
\texttt{Innovate, don't hesitate,} \\
\texttt{Take the leap, don't be late,} \\
\texttt{Create, disrupt, and make it new,} \\
\texttt{That's what innovation's all about, boo.} \\
\\ 
\texttt{Outro:} \\
\texttt{So don't be shy, don't be slow,} \\
\texttt{Innovate, and let your creativity grow,} \\
\texttt{The world needs you, it needs your voice,} \\
\texttt{Make your mark, make some noise, rejoice."} \\
\\ 
\texttt{This song aims to capture the excitement and possibility of innovation, while also emphasizing the importance of taking risks and thinking outside the box. The language is relatable and energetic, with a focus on contemporary references and slang that will resonate with teenagers. The song tells a story about innovation and its impact on teenagers' lives and futures, with a clear and compelling message that encourages listeners to be creative and innovative.} \\
            }
        }
    }
    & & \\ \\

    \theutterance \stepcounter{utterance}  
    & & & \multicolumn{2}{p{0.3\linewidth}}{
        \cellcolor[rgb]{0.9,0.9,0.9}{
            \makecell[{{p{\linewidth}}}]{
                \texttt{\tiny{[GM$|$GM]}}
                {'startend:quotation': {}}
            }
        }
    }
    & & \\ \\

    \theutterance \stepcounter{utterance}  
    & & & \multicolumn{2}{p{0.3\linewidth}}{
        \cellcolor[rgb]{0.9,0.9,0.9}{
            \makecell[{{p{\linewidth}}}]{
                \texttt{\tiny{[GM$|$GM]}}
                \texttt{game\_result = LOSE} \\
            }
        }
    }
    & & \\ \\

    \theutterance \stepcounter{utterance}  
    & & \multicolumn{4}{p{0.6\linewidth}}{
        \cellcolor[rgb]{0.95,0.95,0.95}{
            \makecell[{{p{\linewidth}}}]{
                \texttt{\tiny{[P1$\leftrightarrow$P]}}
                \texttt{Post{-}episode Reflexion (not sent to GM):} \\
\texttt{New plan: My previous attempts at writing a song about innovation were successful in terms of formatting, but they failed to capture the essence of the task, which was to write a song appealing to teenagers. I focused too much on the general theme of innovation and didn't specifically address the tone and language that would resonate with teenagers. However, upon re{-}examining my previous attempts, I realized that I did use some contemporary references and catchy rhymes, but I didn't go far enough in making the language relatable and energetic.} \\
\\ 
\texttt{To improve, I will focus on using more colloquial language, slang, and idioms that are commonly used by teenagers. I will also try to incorporate more specific examples and references to popular culture that are relevant to the teenage audience. Additionally, I will pay closer attention to the rhythm and flow of the song, using more consistent and catchy rhymes to make it more engaging and memorable.} \\
\\ 
\texttt{However, I realize that my previous attempts were too focused on listing off examples of innovation and didn't really tell a story or convey a clear message. To improve, I will try to create a more narrative{-}driven song that tells a story about innovation and its impact on teenagers. I will use the contemporary references and slang to make the story feel more relatable and authentic, and I will focus on creating a clear and compelling message that resonates with the target audience.} \\
\\ 
\texttt{Specifically, I will try to write a song that addresses the following themes:} \\
\\ 
\texttt{* The excitement and possibility of innovation} \\
\texttt{* The importance of taking risks and trying new things} \\
\texttt{* The impact of innovation on teenagers' lives and futures} \\
\texttt{* The need for teenagers to be creative and think outside the box} \\
\\ 
\texttt{To achieve this, I will start by brainstorming a list of contemporary references, slang, and idioms that are commonly used by teenagers. I will then use this list to inform my writing, incorporating these elements into the song in a way that feels natural and authentic. I will also experiment with different rhythms and rhymes, using online tools and resources to help me find the right balance of flow and catchiness.} \\
\\ 
\texttt{To make the song more engaging and memorable, I will try to use a more conversational tone and incorporate storytelling elements, such as metaphors and anecdotes, to make the message more relatable and impactful. I will also try to use a more consistent and catchy rhyme scheme, and experiment with different musical elements, such as melodies and harmonies, to make the song more engaging and memorable.} \\
\\ 
\texttt{Ultimately, my goal is to write a song that not only meets the technical requirements of the task but also resonates with the target audience of teenagers. By focusing on relatable language, catchy rhymes, and a more energetic tone, I am confident that I can create a song that is both engaging and memorable.} \\
            }
        }
    }
    & \\ \\

\end{supertabular}
}

\end{document}
